\documentclass[11pt]{article}
\usepackage{amsmath, amssymb, amsthm, mathtools}
\usepackage[margin=1in]{geometry}
\usepackage{booktabs}

% Shortcuts
\newcommand{\N}{\mathbb{N}}
\newcommand{\Z}{\mathbb{Z}}
\newcommand{\Q}{\mathbb{Q}}
\newcommand{\R}{\mathbb{R}}
\newcommand{\T}{\mathbb{T}}
\newcommand{\set}[1]{\{#1\}}

% Theorem-like environments
\newtheorem{theorem}{Theorem}
\newtheorem{lemma}{Lemma}
\newtheorem{claim}{Claim}
\theoremstyle{definition}
\newtheorem{definition}{Definition}

\title{Problem Set 2}
\author{Michael Duren}
\date{2026-11-05}

\renewcommand{\thesubsection}{\alph{subsection}}

\begin{document}
\maketitle

\section*{Problem 1}

% PROBLEM 1A =================================================================================================================
\subsection{[3 pts]} \textbf{Theorem:} "Every natural number is tall or wide, but not both." (Proof method: proof by cases.)

\begin{proof}
	The proposition:
	\[
		\forall n \in \N, (\text{Tall}(n) \vee \text{Wide}(n)) \wedge \neg(\text{Tall}(n) \wedge \text{Wide}(n))
	\]
	Let $n \in \N$.

	\textbf{Part 1 (existence):} We WTS $\text{Tall}(n) \wedge \text{Wide}(n)$.
		[TODO: show every natural number is tall or wide].

	\textbf{Part 2(uniqueness):} We WTS $\neg(\text{Tall}(n) \wedge \text{Wide}(n))$ by cases on whether $n$ is tall.
	\begin{itemize}
		\item \emph{Case 1: $n$ is tall.} [TODO: show $n$ is not wide.] Then $\neg (\text{Tall}(n) \wedge \text{Wide}(n))$.
		\item \emph{Case 2: $n$ is not tall.} Then $\text{Tall}(n) \wedge \text{Wide}(n)$ is false because the first conjunct fails.
	\end{itemize}
	Combing Parts 1 and 2 gives the theorem.
\end{proof}

\vspace{1cm}
% END PROBLEM =================================================================================================================

% PROBLEM 1B =================================================================================================================
\subsection{[3 pts]} \textbf{Theorem:} "There exists a happy natural number between 10 and 1000 that has
no happy proper divisors" (Proof method to use: choose 25 as the example)

\begin{proof}
	The proposition $\text{P}$:
	\[
		\exists n \in \N, 10 < n < 1000 \wedge \text{isHappy}(n) \wedge (n \text{ has no happy proper divisors})
	\]
	We claim $n = 25$ satisfies the condition.

	\begin{itemize}
		\item \textbf{Range: } $10 < n < 25$ \checkmark
		\item $n$ itself is happy because [TODO: prove]
		\item $n$ has two proper divisors: $1$ and $5$.

		      \begin{itemize}
			      \item WTS that 1 is not a \underline{happy} number because: [TODO: prove why]
			      \item WTS that 5 is not a \underline{happy} number because: [TODO: prove why]
		      \end{itemize}
	\end{itemize}

	Therefore 25 proves the existential claim, proving the theorem.
\end{proof}

\vspace{1cm}
% END PROBLEM =================================================================================================================

% PROBLEM 1C =================================================================================================================
\subsection{[3 pts]} \textbf{Theorem:} “There are no devious natural numbers.” (Proof method to use: regular induction, with base cases 0 and 1.)
\begin{proof}
	The proposition:
	\[
		\forall n \in \N,\ \neg \text{isDevious}(n)
	\]
	Let $P(n) := \neg \text{isDevious}(n)$. We prove $\forall n \in \N,\ P(n)$ by induction.
	\begin{itemize}
		\item \textbf{Base case $P(0)$:} WTS $0$ is not devious. [TODO]
		\item \textbf{Base case $P(1)$:} WTS $1$ is not devious. [TODO]
		\item \textbf{Inductive step:} Let $n \geq 1$, and assume $P(n-1)$ and $P(n)$. WTS $P(n+1)$.
			      [TODO: using that $n-1$ and $n$ are not devious, show $n+1$ is not devious.]
	\end{itemize}
	By induction, $\forall n \in \N,\ \neg \text{isDevious}(n)$.
\end{proof}

\vspace{1cm}
% END PROBLEM =================================================================================================================

\section*{Problem 2}
% RESET SUB SECTION COUNTER 
\setcounter{subsection}{0}

% PROBLEM 2A =================================================================================================================
\subsection{[3 pts]} Befuddled Brynmor tries to help by proving the claim as follows. What is his mistake?
\begin{proof}
	Brynmor substitutes $x_{i-1} \geq 2^{i-1}$ into $-8 x_{i-1}$ as if it produces a lower bound, but since the coefficient is negative the inequality direction flips: $-8 x_{i-1} \leq -8 \cdot 2^{i-1}$. So his conclusion $x_{i+1} \geq 6 \cdot 2^i - 8 \cdot 2^{i-1}$ is unjustified.
\end{proof}
\vspace{1cm}

\vspace{1cm}
% END PROBLEM =================================================================================================================

% PROBLEM 2B =================================================================================================================
\subsection{[3 pts]} Erroneous Erik tries to help by proving the claim as follows. What is his mistake?

\begin{proof}
	The bug is located in Eriks use of his inductive hypothesis $\text{R}(i) := "x_i \geq 3x_{i-1}"$. $\text{R}(i)$ states
	nothing about $x_i > 0$ and consequently doesn't prove this in either the base case or induction step.
	This becomes a circular reference since the goal is to prove for all $i \geq 1, x_i > 0$. If Erik were to
	strengthen his IH, by redefining R as $\text{R}(i) := "x_i \geq 3x_{i-1} \wedge x_i > 0"$ this would prove the
	theorem.
\end{proof}

\vspace{1cm}
% END PROBLEM =================================================================================================================

% PROBLEM 2C =================================================================================================================
\subsection{[3 pts]} Overzealous Zach tries to help by proving the claim as follows. What is his mistake?
\begin{proof}
	Zachs inductive step assumes $\text{S}(i)$ and $\text{S}(i-1)$ but he has only proven one base case.
	The recurrence of $x_{i+1} = 6x_{i} - 8x_{i-1}$ reaches back two steps but he has not proven two base cases.
	If we were to reach back another step from $\text{S}(i)$ we would get $\text{S}(0)$ which is not apart of the original
	sequence.
\end{proof}

\vspace{1cm}
% END PROBLEM =================================================================================================================


% PROBLEM 2D =================================================================================================================
\subsection{[6 pts]} Prove the original claim that for all $i \geq 1, x_i > 0$.

\begin{theorem}
	Let $(x_1, x_2, x_3, \ldots)$ be the sequence of integers defined by $x_1 = 4$, $x_2 = 13$, and
	$x_i = 6 x_{i-1} - 8 x_{i-2}$ for all $i \geq 3$. Then for all $i \geq 1$, $x_i > 0$.
\end{theorem}

\begin{proof}
	We'll use regular induction on the induction hypothesis $\text{IH}(i) := "x_i \geq 3x_{i-1} \wedge x_{i-1} > 0"$ holds for all
	$i \geq 2$.

	\begin{itemize}
		\item \textbf{Base case} $\text{IH}(2)$: $x_2 = 13 \geq 12 = 3x_1$. $x_1 = 4 > 0$ \checkmark
		\item \textbf{Inductive Step}: Suppose $i \geq 2$. Assuming $x_i \geq 3x_{i-1} \wedge x_{i-1} > 0$ (IH$(i)$), WTS that $x_{i+1} \geq 3x_i \wedge x_i > 0$ (IH$(i+1)$).

		      First, from IH$(i)$ we have $x_{i-1} > 0$, so $3x_{i-1} > 0$, and $x_i \geq 3x_{i-1} > 0$. Hence $x_i > 0$, establishing the second conjunct of IH$(i+1)$.

		      If $x_i \geq 3x_{i-1}$, we can rewrite this inequality by multiplying both sides by $\frac{8}{3}$ to get $\frac{8}{3}x_i \geq 8x_{i-1}$. Rewriting the recurrence $x_{i+1} = 6x_i - 8x_{i-1}$ as $x_{i+1} = \frac{10}{3}x_i + \frac{8}{3}x_i - 8x_{i-1}$, and noting that $\frac{8}{3}x_i - 8x_{i-1} \geq 0$ from above, we get $x_{i+1} \geq \frac{10}{3}x_i$. Since $x_i > 0$, $\frac{10}{3}x_i \geq 3x_i$, so $x_{i+1} \geq 3x_i$.

		      Combined with $x_i > 0$, we have $x_{i+1} \geq 3x_i > 0$, establishing both conjuncts of IH$(i+1)$. This completes the inductive step for all $i \geq 2$.
		\item \textbf{Conclusion:} By induction, $\text{IH}(i)$ holds for all $i \geq 2$. Since $\text{IH}(i)$ shows that $x_{i-1} > 0$ taking
		      all of $i = 2, 3, 4 \ldots$ gives $x_1 > 0, x_2 > 0, x_3 > 0 \ldots$ covering every $x_i$ for $i \geq 1$. Therefore $x_i > 0$ for all of $i \geq 1$.
	\end{itemize}

\end{proof}

\vspace{1cm}
% END PROBLEM =================================================================================================================

\section*{Problem 3}
% RESET SUB SECTION COUNTER 
\setcounter{subsection}{0}

% PROBLEM 3A =================================================================================================================
\subsection{[2 pts]}  Compute the full sequence of transitions from starting state (5, 22, 0). Does it
correctly compute 5 · 22 = 110?.

\begin{tabular}{cccc}
	\toprule
	r   & s  & a   & rule        \\
	\midrule
	5   & 22 & 0   & \text{even} \\
	10  & 11 & 0   & \text{odd}  \\
	20  & 5  & 10  & \text{odd}  \\
	40  & 2  & 30  & \text{even} \\
	80  & 1  & 30  & \text{odd}  \\
	160 & 0  & 110 & \text{done}
	\bottomrule
\end{tabular}

\vspace{1cm}
% END PROBLEM =================================================================================================================

% PROBLEM 3B =================================================================================================================
\subsection{[2 pts]}
\textbf{State machine rules}:
The states are triples $(r, s, a) \in \N^3$, with start state $(x, y, 0)$.
Transitions are defined for $s > 0$:
\[
	(r, s, a) \mapsto
	\begin{cases}
		(2r,\ s/2,\ a)         & \text{if } (s \text{ is even} \wedge  s > 0) \\
		(2r,\ (s-1)/2,\ a + r) & \text{if } (s \text{ is odd} \wedge s > 0)   \\
	\end{cases}
\]

\textbf{\emph{Invariant.}} Define $P(r,s,a) \coloneqq rs + a = xy$
\\
\begin{theorem*}
	Theorem: $P$ is an invariant of the state machine.
\end{theorem*}

\begin{proof}
	By \emph{invariant} principle.
	\begin{itemize}
		\item $P(r, s, a)$ holds at the start state $(x, y, 0)$ because
		      $rs + a = xy + 0 = xy$.
		\item $P(r, s, a)$ is preserved: for any transition
		      $(r, s, a) \to (r', s', a')$ with $P(r, s, a)$ true,
		      we show $P(r', s', a')$ holds.
		      Since $s \in \mathbb{N}$, if $s = 0$ there is no transition,
		      so we consider the cases where $s > 0$.
		      \begin{itemize}
			      \item \textbf{Case 1: $s$ is even.}
			            \\
			            The next state is $(r', s', a') = (2r,\; s/2,\; a)$.
			            We know $xy = rs + a$ by assumption $P(r, s, a)$,
			            and WTS $r's' + a' = xy$.
			            We have $r's' + a' = 2r(s/2) + a = rs + a = xy$.
			            Thus $P(r', s', a')$ holds.
			            \\
			      \item \textbf{Case 2: $s$ is odd.}
			            \\
			            The next state is $(r', s', a') = (2r,\; (s-1)/2,\; a + r)$.
			            We know $xy = rs + a$ by assumption $P(r, s, a)$,
			            and WTS $r's' + a' = xy$.
			            We have $r's' + a' = 2r \cdot \frac{s-1}{2} + a + r
				            = r(s-1) + a + r = rs - r + r + a = rs + a = xy$.
			            Thus $P(r', s', a')$ holds.
		      \end{itemize}
		\item \textbf{Conclusion:} By the invariant principle,
		      $P(r, s, a)$ holds for all reachable states.
	\end{itemize}
\end{proof}

\vspace{1cm}
% END PROBLEM =================================================================================================================

\end{document}
